%=============================================================================
%  04_architecture.tex  -  Chapter 3: System Architecture
%=============================================================================
\chapter{System Architecture}
\label{ch:arch}

\section{Overview}
Universal Analyst is a two-tier application: a React single-page front end and a
FastAPI backend, with the backend orchestrating a set of AI and data
subsystems and delegating language generation to a local Ollama runtime. The
front end never talks to Ollama directly; every capability is mediated by the
backend so that validation, grounding, and the no-fabrication rule can be
enforced in one place.

\begin{figure}[h]
\centering
\begin{tikzpicture}[node distance=8mm and 12mm]
  % Frontend
  \node[ua-fe,text width=3.2cm] (ui) {Frontend\\\footnotesize React + Vite + Tailwind};

  % Backend layer
  \node[ua-be,below=14mm of ui,text width=3.6cm] (routes) {API Layer (FastAPI)\\\footnotesize /upload /chat /visualize\\/skills /sessions /clean /predict};
  \node[ua-be,left=of routes,text width=2.3cm] (router) {Intent\\Router};
  \node[ua-be,right=of routes,text width=2.3cm] (sess) {Session\\Manager};

  % AI/data layer
  \node[ua-ai,below=16mm of router,text width=2.6cm] (data) {Ingestion\\Profiler\\Auto-EDA};
  \node[ua-ai,below=16mm of routes,text width=2.6cm] (rag) {RAG\\\footnotesize ChromaDB + MiniLM};
  \node[ua-ai,below=16mm of sess,text width=2.6cm] (skills) {Skills\\\footnotesize clean/predict/report/viz};

  \node[ua-ai,below=12mm of rag,text width=3.0cm] (llm) {Ollama Client\\\footnotesize Gemma 4 (12B / e4b)};

  % Output layer
  \node[ua-out,below=12mm of data,text width=2.6cm] (viz) {Plotly\\Chart Factory};
  \node[ua-out,below=12mm of skills,text width=2.6cm] (pdf) {ReportLab\\+ Kaleido PDF};

  % edges
  \draw[ua-arrow] (ui) -- node[right,font=\scriptsize,text=uaSlate]{/api/*} (routes);
  \draw[ua-arrow] (routes) -- (router);
  \draw[ua-arrow] (routes) -- (sess);
  \draw[ua-arrow] (router) -- (data);
  \draw[ua-arrow] (routes) -- (rag);
  \draw[ua-arrow] (sess) -- (skills);
  \draw[ua-arrow] (data) -- (rag);
  \draw[ua-arrow] (rag) -- (llm);
  \draw[ua-arrow] (skills) -- (llm);
  \draw[ua-arrow] (skills) -- (viz);
  \draw[ua-arrow] (skills) -- (pdf);
  \draw[ua-arrowl] (llm) to[bend left=12] node[left,font=\scriptsize,text=uaAmber]{tokens / plans} (routes);

  % layer labels (anchored to the left-column nodes)
  \node[font=\scriptsize\bfseries,text=uaIndigo,rotate=90] at ([xshift=-7mm]ui.west) {CLIENT};
  \node[font=\scriptsize\bfseries,text=uaGreen,rotate=90] at ([xshift=-7mm]router.west) {BACKEND};
  \node[font=\scriptsize\bfseries,text=uaAmber,rotate=90] at ([xshift=-7mm]data.west) {AI / DATA};
  \node[font=\scriptsize\bfseries,text=uaIndigoD,rotate=90] at ([xshift=-7mm]viz.west) {OUTPUT};
\end{tikzpicture}
\caption{High-level component architecture. The backend is the single point at
which grounding, validation, and the no-fabrication rule are enforced.}
\label{fig:arch}
\end{figure}

\section{Component responsibilities}
\begin{table}[h]
\centering
\caption{Backend packages and their responsibilities.}
\small
\begin{tabularx}{\linewidth}{@{}L{3.2cm} X@{}}
\toprule
\textbf{Package} & \textbf{Responsibility} \\
\midrule
\code{app.api} & HTTP route handlers: chat, upload, sessions, visualizations,
skills, clean, predict, health. \\
\code{app.core} & Ingestion, chunking, profiling, and the in-memory session
manager. \\
\code{app.eda} & Automated exploratory data analysis over the profile. \\
\code{app.llm} & Ollama HTTP client, prompt templates, and token counting. \\
\code{app.rag} & ChromaDB store, sentence-transformers embedder, retriever, and
context builder. \\
\code{app.skills} & Cleaning, prediction, report, and report-narrative
implementations. \\
\code{app.skills\_registry} & Markdown definitions that steer the model per
skill. \\
\code{app.visualization} & Plotly chart factory, chart-type logic, and image
export. \\
\code{app.models} & Pydantic schemas and enumerations shared across the app. \\
\bottomrule
\end{tabularx}
\end{table}

\section{Request lifecycle}
Figure~\ref{fig:lifecycle} traces a full session: an upload that profiles and
indexes the data, a question that is grounded via retrieval and streamed from
the model, an optional chart that is validated before rendering, and a report
export.

\begin{figure}[h]
\centering
\begin{tikzpicture}[
  actor/.style={font=\footnotesize\bfseries,text=uaInk},
  lifeline/.style={uaLine,thick},
  msg/.style={-{Stealth[length=2.4mm]},uaSlate,thick,font=\scriptsize},
  ret/.style={-{Stealth[length=2.4mm]},uaAmber,thick,dashed,font=\scriptsize},
  x=2.05cm,y=-0.62cm]

  \foreach \i/\name in {0/Analyst,1/{React UI},2/FastAPI,3/{Ingest\,/\,Profiler},4/{RAG},5/{Gemma\,4},6/{Plotly\,/\,PDF}}{
    \node[actor] (h\i) at (\i,0) {\name};
    \draw[lifeline] (\i,-0.6) -- (\i,-19.5);
  }

  \draw[msg] (0,-1.6) -- node[above]{upload} (1,-1.6);
  \draw[msg] (1,-2.4) -- node[above]{POST /upload} (2,-2.4);
  \draw[msg] (2,-3.2) -- node[above]{parse + profile + EDA} (3,-3.2);
  \draw[msg] (3,-4.0) -- node[above]{chunk + embed} (4,-4.0);
  \draw[ret] (2,-4.9) -- node[below]{session\_id, profile, eda} (1,-4.9);

  \draw[msg] (0,-6.1) -- node[above]{ask question} (1,-6.1);
  \draw[msg] (1,-6.9) -- node[above]{POST /chat (stream)} (2,-6.9);
  \draw[msg] (2,-7.7) -- node[above]{retrieve top-k} (4,-7.7);
  \draw[ret] (4,-8.5) -- node[below]{grounded chunks} (2,-8.5);
  \draw[msg] (2,-9.3) -- node[above]{prompt (+images?)} (5,-9.3);
  \draw[ret] (5,-10.1) -- node[below]{thinking / tokens / plan} (2,-10.1);

  \node[font=\scriptsize\itshape,text=uaAmber] at (3.5,-11.0) {alt [chart requested]};
  \draw[msg] (2,-11.7) -- node[above]{validate plan} (6,-11.7);
  \draw[ret] (6,-12.5) -- node[below]{chart event} (1,-12.5);

  \draw[ret] (2,-13.6) -- node[below]{tokens / usage / done} (1,-13.6);

  \draw[msg] (0,-14.9) -- node[above]{export report} (1,-14.9);
  \draw[msg] (1,-15.7) -- node[above]{POST /skills/report} (2,-15.7);
  \draw[msg] (2,-16.5) -- node[above]{build PDF} (6,-16.5);
  \draw[ret] (6,-17.3) -- node[below]{download\_report} (1,-17.3);
\end{tikzpicture}
\caption{End-to-end request lifecycle, from upload to grounded answer to report
export. Dashed amber arrows are streamed or returned responses.}
\label{fig:lifecycle}
\end{figure}

\section{Routing and paths}
The backend serves its routes at the root of port \code{8000}. The front end,
however, always calls paths under \code{/api/*}. In every deployment mode a
proxy --- the Vite dev server in development, nginx in production, and a small
forwarding server in the notebook --- rewrites \code{/api/*} onto the backend
while stripping the \code{/api} prefix. This keeps the front-end code identical
across environments and avoids hard-coding backend hostnames.

\begin{uanote}[Why a single mediating backend]
Because grounding (retrieval), validation (chart plans), token accounting, and
error honesty all live behind one API, the front end can stay thin and the
guarantees stay consistent. There is exactly one place where a chart can be
rejected or an error surfaced --- not several.
\end{uanote}

\section{Deployment topology}
Three runtime shapes are supported from the same codebase, summarized in
Table~\ref{tab:topology} and detailed in Chapter~\ref{ch:deploy}.

\begin{table}[h]
\centering
\caption{Supported deployment topologies.}
\label{tab:topology}
\small
\begin{tabularx}{\linewidth}{@{}L{2.7cm} L{3.2cm} X@{}}
\toprule
\textbf{Mode} & \textbf{Front end} & \textbf{Notes} \\
\midrule
Notebook (cloud) & Forwarded Vite dev server & \code{main.ipynb} installs deps,
pulls the Gemma tags, and starts both services on a GPU box. \\
Docker (dev) & Vite dev server in a Node image & Bind-mounts source; Ollama runs
on the host via \code{host.docker.internal}. \\
Docker (prod) & Built static assets behind nginx & \code{docker-compose.prod.yml}
serves optimized images. \\
\bottomrule
\end{tabularx}
\end{table}
